\documentclass[11pt]{article}
\usepackage[UTF8]{ctex}
\usepackage[a4paper,margin=1in]{geometry}
\usepackage{booktabs}
\usepackage{array}
\usepackage{hyperref}
\usepackage{xcolor}
\usepackage{setspace}
\usepackage{longtable}
\usepackage{amsmath}
\setstretch{1.15}
\hypersetup{colorlinks=true,linkcolor=blue,citecolor=blue,urlcolor=blue}

\title{多请求常驻分叉下的 vLLM Q16 同核 Attention 复用与容量分析：\\
Qwen3.5-35B 4K 场景的可复现审计}
\author{匿名}
\date{\today}

\begin{document}
\maketitle

\begin{abstract}
本文报告 Qwen3.5-35B-A3B 上的一个受控可复现实验：在单批量（batch size=1）、Q16 精度、single stream、单 GPU 进程内，比较两种 same-kernel 文档所有权策略在常驻请求数
$N=\{1,2,4,8,16,32\}$ 下的解析容量与一致性。我们在 PG-19 train-only 的同一文档集合上，使用同一份约 4095-token 的预填充上下文（worst-case partial-tail 情况）、相同几何/硬件/kernel 条件，得到两点结论：第一，两条 arm 在 token、full-vocab logits、最终 logical K/V 与 GDN state 上达到逐步一致（token-level + full-logit exact）。第二，可复用文档策略的解析账本规模随 $N$ 近似线性下降，新增一个常驻请求约节省 85 MiB，其中 80 MiB 来自避免物理文档整块复制；该优势对应的 PyTorch allocator absolute peak allocated 中位数在 $N=32$ 时下降约 2.661 GiB（约占 fresh 的 3.83\%）。  
为避免范围漂移，本文仅讨论 \textbf{same-kernel same-batch 条件下的解析账本与 correctness 门禁}，不涉及吞吐/延迟、scheduler、ragged 批处理、多文档/多源泛化、Q8/Q4、下游任务指标。

\textbf{贡献结构。}
\begin{itemize}
  \item 在单流、同核、single-batch 条件下，给出 $N\in\{1,2,4,8,16,32\}$ 的多常驻实验与证据闭环。
  \item 验证同核复用与 fresh baseline 在 token、full-vocab、logical K/V、GDN state 上的 exact。
  \item 给出可复验的解析容量线性模型（fresh 与 reuse）与 allocator 中位数差异。
\end{itemize}
本文不声明并发吞吐、延迟加速、NVML 总显存提升或多文档泛化，也不声明与 HF-eager 的 logits bitwise 等价。
\end{abstract}

\section{引言}
\textbf{问题定义。} 在长上下文推理里，常驻请求（resident requests）数目增长时，若每个请求都拷贝完整文档，解析内存会快速线性膨胀；但若共享文档并隔离请求私有区，是否能在不改变生成结果的前提下降低解析内存，是本工作关注的核心问题。  
\textbf{已有差距。} 前序实验已验证了 single request 时的 same-kernel exact，但没有系统扫描常驻请求数量 $N>1$ 的容量-一致性权衡。  
\textbf{本文贡献与边界。} 本文给出 one-kernel、batch1、Q16、PG-19 train-only 下的同一文档同源跨 $N$ 可比实验，提供完整证据包、审计门禁和边界。  
\textbf{结论边界预告。} 我们只证明“解析账本改善与 correctness 保持”，并不推广到吞吐、scheduler、ragged batch、多文档、Q8/Q4 或下游任务 F1/EM 改善。本文也不主张 timing/TTFT/TPOT speedup。

\section{方法与实验设置}
\textbf{总体思路。} 我们把单请求实验扩展为同一 4K 文档的常驻多请求曲线，固定一套几何与数据治理，同时比较两种文档所有权策略。  
\textbf{主命题。} 在不改变输出结果的前提下，量化共享文档池方案对“常驻请求规模”是否带来可复用容量优势。
\textbf{对比定义。}
\begin{itemize}
  \item \textbf{fresh full-copy control}：在同一 common prefill/pack 后，为每个请求分配独立 document+private 区；每个请求复制完整物理文档块。  
  \item \textbf{shared-document reuse}：文档块共享为只读 source arena；每请求仅持有私有页面用于 partial-tail COW 与 append。  
\end{itemize}

\textbf{模型与内核。} 使用 Qwen3.5-35B-A3B，40 层 Transformer，其中 full-attention 层索引为 3,7,11,15,19,23,27,31,35,39。内核为 vLLM 0.26 unified attention，符号为 \verb|triton_unified_attention.unified_attention|，Q16，page size 128，batch 1。环境固定 \texttt{PyTorch 2.11.0+cu129}, \texttt{Transformers 5.14.1}, \texttt{Triton 3.6.0}, \texttt{FlashInfer 0.6.14}。
\[
N\in\{1,2,4,8,16,32\},\quad \text{query}=32,\quad \text{gen}=8,\quad \text{doc}=4095
\]

\textbf{数据与治理。} 每个 rank 使用不同 PG-19 train book window；每个 book 固定冻结 32 条互不重叠 32-token query（stride=64），无 synthetic marker，无 source6-9/68-99/test-v2 消耗，保证与生产日志链路一致。文档长度为 4095（$mod\,128=127$），用于偏压 tail-copy 的 worst-case 应力场。  
\textbf{常驻规模。} 每个 rank 全量运行 $N=\{1,2,4,8,16,32\}$；即每个 rank 对所有 N 点都有同源 query 前缀可比（即 N 嵌套可比）。  
\textbf{执行顺序与资源。} 8 张 H20-3e（单试验 8 GPU）；每个 rank 顺序处理不同 N，但每个 rank 内 1 个 source + 2/4/8/16/32 个常驻请求对象在 setup 与生成阶段并发驻留。单流执行采用 round-major 顺序（先各请求做 round 0，再各请求做 round 1，依次到 round 7），不构成 scheduler 并发。

\section{正确性与一致性}
\textbf{一致性约束。} 每个 raw shard 先记录 token、每步 full-vocab logit、logical K/V、GDN state，并在聚合时重新计算并对照，拒绝仅信任顶层布尔量。
\textbf{核心正确性结果。}
\begin{itemize}
  \item fresh 与 reuse 在 8 个 rank、全部 N、全部请求、每个生成 token 的 full-vocab 与 token 均 exact（比例 1.0）。  
  \item 两个 arm 的 final logical K/V 与 60-tensor GDN state exact。  
  \item 10 个 attention 层共用同一 unified\_attention descriptor；dense fallback 与 full concat 的计数为 0。  
  \item cross-$N$ prefix isolation：同 rank 同一 request i 在不同 N 的输出在同一 seed/seedless 前缀下 exact，避免“共享状态污染两 arm 同时污染”的假阳性。  
\end{itemize}
\textbf{未能支持的外延。} 本实验不证明 HF-eager 与 vLLM logits/everything exact；HF 对照仅作为 non-blocking 诊断，不能替代 same-kernel 主结论。

\section{容量与资源结果}
\textbf{解析容量解析。} 单层 physical page block 为 2,621,440B（十层为 26,214,400B）。在 4095 token 下，document allocation 为 32 pages，即 80MiB；有效 payload 80MiB 减去少量 padding；每请求 private reservation 为 2 pages（约 5MiB）。  
\textbf{公式与解释。}
\begin{align}
\text{fresh pool}(N) &= 80\mathrm{MiB} + 90N\mathrm{MiB}\\
\text{reuse pool}(N) &= 80\mathrm{MiB} + 5N\mathrm{MiB}\\
\text{savings}(N) &= 85N\mathrm{MiB}
\end{align}
其中 80MiB 来自避免的文档物理拷贝，5MiB 为每请求私有页保留，且每个请求在 active 阶段仍保留 partial-tail/append。  
\textbf{容量曲线（10 full-attn 层，8 rank 媒体中位数）：}
\begin{table}[h]
\centering
\begin{tabular}{cccccc}
\toprule
N & fresh (MiB) & reuse (MiB) & controlled savings (MiB) & avoided doc copy (MiB)\\
\midrule
1 & 170 & 85 & 85 & 80\\
2 & 260 & 90 & 170 & 160\\
4 & 440 & 100 & 340 & 320\\
8 & 800 & 120 & 680 & 640\\
16 & 1520 & 160 & 1360 & 1280\\
32 & 2960 & 240 & 2720 & 2560\\
\bottomrule
\end{tabular}
\caption{解析容量中位数结果，来自 raw 重放回归后 fit。}
\end{table}
\textbf{说明。} 表里是解析账本量，不是模型总内存，也不是 NVML 总显存。每个结果均来自 raw shard 重放并带层级一致性校验。

\section{分配器与指标解释}
\textbf{统计口径。} 我们使用 PyTorch allocator 的 allocated/reserved 指标，不宣称与系统级 NVML peak 等价。  
\textbf{关键差值。} 在 $N=32$ 下，absolute production peak allocated 为 fresh 74,623,183,360B、reuse 71,765,915,136B，差值 2,857,268,224B（约 2.661 GiB）。  
\textbf{为什么不算时延提升。} 仅有单流 round-major 顺序（不是 ABBA，不是 scheduler 或 continuous batching）和 single-GPU rank 内存测量，故不报告 TTFT/TPOT/吞吐 speedup。

\begin{table}[h]
\centering
\begin{tabular}{ccccc}
\toprule
N & fresh setup+gen current delta & reuse setup+gen current delta & 差值 & fresh/reuse allocated peak (GiB)\\
\midrule
1 & 147.135 MiB & 61.883 MiB & 85.252 MiB & 64.896 / 64.813\\
2 & 298.522 MiB & 126.766 MiB & 171.757 MiB & 65.048 / 64.881\\
4 & 591.286 MiB & 249.520 MiB & 341.767 MiB & 65.345 / 65.011\\
8 & 1177.831 MiB & 495.045 MiB & 682.786 MiB & 65.936 / 65.270\\
16 & 2355.406 MiB & 990.081 MiB & 1365.325 MiB & 67.126 / 65.792\\
32 & 4705.065 MiB & 1980.162 MiB & 2724.903 MiB & 69.498 / 66.837\\
\bottomrule
\end{tabular}
\caption{8 rank 中位数；reserved 主要受 allocator caching 行为影响，本文不转化为可移植容量公式。}
\end{table}

\section{可复现性与治理}
\textbf{身份与审计。} 运行为 Job 237580 / Trial 1840837（Complete）。阶段顺序为
\texttt{00\_start $\rightarrow$ 01\_static\_preflight\_ok $\rightarrow$ 02\_resident\_shards\_ok $\rightarrow$ 99\_done}，无 FAILED。  
\textbf{验证闭环。} code/model/windows/manifest、PG-19 4095 windows、query bank、raw shards、ledger 均逐项 hash 复核；aggregate 不信任 shard 顶层布尔，而是重放 raw call、trajectory、allocator、storage 与 query provenance。  
\textbf{证据链摘要。} 关键 SHA：summary \texttt{cbb5d861...f671c0401}，scientific ledger \texttt{14fd1c0c...faf363}（35/35 OK）。

\section{讨论与局限}
\textbf{可支持的结论。} 在本文条件下，可复现实验表明 reuse 策略在保持 same-kernel exact 的同时，显著减少解析扩展，且可线性外推到 N=32。  
\textbf{不应扩展的结论。} 本实验不证明总显存缩放 12.3x；不证明多文档场景、Q8/Q4、ragged-batch/continuous batching、scheduler 并发或下游任务指标。  
\textbf{后续方向。} 下一步可在不改变 correctness 门禁的前提，加入：
\begin{itemize}
  \item 与 longbench/quality 的弱关联任务映射（避免 overclaim）；
  \item 与生产 scheduler 的 decoupled timing 评估；
  \item 不同文档长度与尾部分布下的 worst-case/average-case 曲线分离。  
\end{itemize}

\section{结论}
\textbf{核心结论。} 本文在严格 same-kernel same-batch 条件下，给出 Q16 常驻请求复用的可复现实验：$N$ 增大时解析账本规模按模型可解释公式线性下降，并维持 token/logit/KV/GDN exact。  
\textbf{实用读数。} 到 $N=32$，单纯解析池层面可节省 $2720$ MiB（$85\times32$），且 allocator absolute peak allocated 中位数下降约 $2.661$ GiB；该值仍受缓存行为与 PyTorch 分配语义影响。  
\textbf{严格边界。} 本工作并不推广并发吞吐，亦不把上述账本节省直接写成全模型、NVML、或全系统级显存改善。
\end{document}
